\section{A common power gain in the envelope}\label{sec:envelope}

We construct \(h=h_k\) with a positive test \eqref{eq:test} wherever
packets will be centered, while making the integral of its local
oscillation scale grow by a power of \(k\).

\begin{proposition}\label{prop:envelope}
There are a small \(C^1\) neighborhood \(\mathcal G\) of \(g_0\),
fixed regions \(V\Subset\Omega\subset\{t>1/2\}\), and common
\(a_*,c,\gamma>0\), such that every fixed smooth \(G\in\mathcal G\)
admits envelopes \(h_k\) satisfying
\begin{equation}\label{eq:envelope}
\begin{gathered}
 h_k=-a_*t+c_k\quad(t\le1/2),\qquad
 \|h_k\|_{C^j}\le C_j k^{\delta(j+1)}\quad(j\ge0),\\
 \max_e J_{h_k}(e)\ge c k^{-\delta}\quad(t\ge1/4),\qquad
 \int_V(1+|\nabla_Gh_k|_G^2)^{1/2}\,\dd x\ge c k^\gamma .
\end{gathered}
\end{equation}
The additive constants \(c_k\) may be any prescribed bounded sequence.
For each fixed derivative order the constant and frequency threshold
can depend on that order and on the fixed smooth metric.
The exponent \(\gamma\) is independent of these orders, all later
residual accuracies, and all round gluing collars.
\end{proposition}

\subsection{A seed and a repeated local modification}

The round seed \(\widetilde h=-a_*t\) obeys
\[
 |\nabla_{g_0}\widetilde h|^2=4a_*^2t(1-t),\qquad
 \nabla_{g_0}^2\widetilde h=2a_*(tg_0-\dd y\cdot\dd y).
\]
Since \(\dd y:T_xB\to\R^3\) has a four-dimensional domain, there
is a unit vector in its kernel. It is an admissible test, with Hessian
value \(2a_*t\). The Hessian norm is at most \(2a_*\), so for
\(t\ge1/4\) and \(a_*^2\le1/8\),
\[
 J_{\widetilde h}(e)\ge a_*/2-2a_*^3\ge a_*/4.
\]
At a point with \(t=1\), the Hessian has two eigenvalues \(2a_*\)
and two zero eigenvalues. Add a sufficiently small smooth quadratic
bump there, supported in \(\{t>1/2\}\), to make the Hessian positive
definite at that point. Throughout the bump support two positive
directions remain and the gradient stays sufficiently small. Their
two-plane intersects \(\ker\dd h\), leaving a positive test even at
a critical point. Away from that support the old kernel test applies.
Choose a nearby region \(U_0\) where the Hessian is positive definite
and the gradient is separated from zero. These properties persist in
a common small \(C^1\) metric neighborhood: the Hessian of a fixed
function changes by the connection difference contracted with its
differential. Fix \(U_0\Subset V\Subset\Omega\).

A cube \(P\) of side \(\nu S\), \(S\le1\), is called prepared for \(h\)
if, on its closure,
\begin{equation}\label{eq:prepared}
\begin{gathered}
 |\nabla h|>0,\qquad
 |\nabla^2h|\le C_{\rm p}|\nabla h|/S,\\
 \text{there is a two-plane }E\subset\ker\dd h
 \text{ with }J_h(e)\ge c_{\rm p}|\nabla h|/S
 \text{ for every unit }e\in E.
\end{gathered}
\end{equation}
The plane need not depend continuously on the point. The seed is
prepared at scale \(1\) on a sufficiently small cube in \(U_0\),
for fixed large \(C_{\rm p}\) and small \(c_{\rm p}>0\).

\begin{lemma}\label{lem:corrugation}
One can fix a reciprocal integer \(\nu>0\), constants
\(0<\alpha\le1\le\beta\), and \(\zeta>0\), uniformly under the
common metric bounds, so that every prepared parent \(P\) admits a
smooth modification \(v=h+w\), supported in its interior, with
\[
 \alpha|\nabla h|\le|\nabla v|\le\beta|\nabla h|,\qquad
 \max_eJ_v(e)\ge c|\nabla h|/S.
\]
Put \(\ell=\nu^2S\). Some cubes of side \(\nu\ell\) inside \(P\)
are prepared at scale \(\ell\), with the same constants, and their
union \(\mathcal R_P\) has
\begin{equation}\label{eq:gain}
 \int_{\mathcal R_P}|\nabla v|\,\dd x
 \ge(1+\zeta)\int_P|\nabla h|\,\dd x.
\end{equation}
If \(s=|\nabla h(x_P)|\) at the parent center, then
\begin{equation}\label{eq:increment}
 \|D^jw\|_\infty\le C_j s\ell^{1-j}\qquad(j\ge0).
\end{equation}
Neither the modification nor these bounds use higher derivatives
of the input \(h\).
\end{lemma}

\begin{proof}
Choose a smooth \(F:[0,1]\to\R\), positive inside, equal to its
argument near zero and flat at one. In a finite collection of
disjoint balls of radii \(R_i\) in a unit three-cube, define
\[
 \psi(r)=-\int_r^{R_i}F(q/R_i)\,\dd q,\qquad
 f=F(r/R_i),\qquad u_0=f/r.
\]
Extend by zero and periodically. At the center \(\psi\) is a constant
plus \(r^2/(2R_i)\), so it is smooth there; flatness makes the boundary
extension smooth. With radial unit vector \(\sigma\),
\begin{equation}\label{eq:radial}
 D\psi=f\sigma,\qquad
 D^2\psi=f'\sigma\otimes\sigma+u_0(I-\sigma\otimes\sigma).
\end{equation}
The inequalities \(f^2|f'|\le C u_0\) and \(f^2\le C u_0\) hold.
The first constant does not depend on a ball radius, since
\(f^2|f'|/u_0=vF(v)|F'(v)|\), \(v=r/R_i\).
Also \(u_0=0\) implies \(f=f'=\psi=0\).
At a center the radial Hessian is independent of the choice of \(\sigma\).

Fix a small amplitude \(a>0\) using that radius-independent bound.
The mean gradient factor inside any ball is
\[
 m(a)=3\int_0^1v^2\sqrt{1+a^2F(v)^2}\,\dd v>1.
\]
Let \(\eta_0=(m(a)-1)/8\). Choose the finite packing to have volume
fraction \(\varpi\) with \(\varpi m(a)>1+6\eta_0\). Such finite
packings can approach full volume: finitely many grid cubes capture
a fixed fraction of any remaining open volume; smaller inscribed
balls capture a fixed fraction of those cubes. Repeating this
finitely often leaves an arbitrarily small uncovered volume and
keeps closures disjoint. Reduce each ball to a concentric core
of relative radius \(1-\eta_{\rm c}\), so that the integral per
period over the cores is \(C_{\rm core}>1+5\eta_0\).
Choose an inner-parent volume fraction \(q_{\rm in}<1\) with
\(q_{\rm in}C_{\rm core}>1+4\eta_0\).
A smooth parent cutoff \(\chi\), equal to one on that inner cube
and flat on its zero set, is now fixed.

Take affine coordinates \(z=(z_1,z')\) that make the metric
Euclidean at \(x_P\) and \(\nabla h(x_P)=s\partial_{z_1}\).
The preparation bound and the common connection bound give
\(|D(\dd h)|\le C|\dd h|/S\). Integration along parent segments
of length \(O(\nu S)\) gives
\begin{equation}\label{eq:parent}
 \dd h/s=\dd z_1+O(\nu),\qquad G_{ij}=I+O(\nu),
 \qquad |\nabla h|/s=1+O(\nu).
\end{equation}
Rescale \(\chi\) to the parent and set
\begin{equation}\label{eq:modify}
 v=h+s\ell a\chi\psi(z'/\ell).
\end{equation}
The ratio between the pattern scale and cutoff scale is \(\nu\);
hence
\(\dd v/s=\dd z_1+a\chi f\sigma\cdot\dd z'+O(\nu)\).
This proves the gradient comparison.

We give the Hessian check also where the favorable coefficient
vanishes. Put \(D=\nabla v/\sqrt{1+|\nabla v|^2}\) and
\[
 \mathcal F=\{e:\dd v(e)=0,\ \sigma\cdot e'=0\}.
\]
This is a two-plane. For metric-unit \(e\in\mathcal F\),
\[
 |e'|^2\ge1/2,\qquad |\sigma\cdot D'|\le C_0a\chi f+C\nu.
\]
After dividing the new Hessian test by \(s/\ell\), its principal
contribution from \eqref{eq:radial} is
\[
 a\chi\{f'(\sigma\cdot D')^2+
            u_0(|\Pi_{\sigma^\perp}D'|^2+|e'|^2)\}.
\]
Choose \(a\) so that
\(2C_0^2a^2\sup_v vF(v)|F'(v)|\le1/4\).
The displayed expression is then at least
\(a\chi u_0/4-Ca\nu^2\chi|f'|\).
Mixed cutoff terms on \(e,e\) vanish. On \(D,D\) they are bounded by
\[
 C\nu f|D_{\rm sc}\chi|(\chi f+\nu).
\]
The remaining cutoff and connection errors are bounded by
\[
 C\nu^2|\psi||D_{\rm sc}^2\chi|
       +C\ell(\chi f+\nu).
\]
Here \(D_{\rm sc}\) denotes differentiation at the parent cutoff scale.

Let \(X=\chi u_0\). On the compact set of cutoff and period positions
and bounded affine frames, every continuous \(b\ge0\) that vanishes
when \(X=0\) satisfies
\[
 \sup\frac{\nu^2b}{X+\nu^2}\longrightarrow0.
\]
For \(X\ge\tau\), bound the ratio by \(\nu^2\sup b/\tau\);
for \(X<\tau\), bound it by \(\sup_{X<\tau}b\), then let
\(\tau\downarrow0\). Apply this to
\(\chi|f'|,\ f|D_{\rm sc}\chi|,\ |\psi||D_{\rm sc}^2\chi|\);
their vanishing follows from flatness and \(u_0=0\Rightarrow f=f'=\psi=0\).
The other mixed term is \(O(\nu X)\), and
\(\chi f\le C\sqrt X\), \(\ell\le\nu^2\) imply
\(\ell(\chi f+\nu)/(X+\nu^2)=O(\nu)\).
Consequently, uniformly over admissible parents,
\begin{equation}\label{eq:newhessian}
 \frac{\ell}{s}\{\nabla^2(v-h)(D,D)+\nabla^2(v-h)(e,e)\}
 \ge c_aX-\varepsilon(\nu)(X+\nu^2),\qquad
 \varepsilon(\nu)\to0.
\end{equation}

The old Hessian is \(O(\nu^2)\) in these normalized units.
If \(X/\nu^2\to\infty\), the new positive term dominates it.
If this ratio stays bounded, the old and new gradients differ by a
relative \(O(\nu)\). The map \(T(p)=p/\sqrt{1+|p|^2}\) obeys
\[
 |T(q)-T(p)|\le2\epsilon
 \quad\text{when }|q-p|\le\epsilon|p|,\quad\epsilon\le1/2.
\]
This follows by integrating its derivative along a segment whose
distance from zero is at least \(|p|/2\). A rotation \(R=I+O(\nu)\)
takes \(\ker\dd h\) to \(\ker\dd v\). Inside this three-dimensional
space, \(RE\) and \(\mathcal F\) have intersection dimension at least
\(2+2-3=1\). In that direction the old positive test changes by
\(o(\nu^2)\); the new contribution in \eqref{eq:newhessian} is
bounded below by \(-o(\nu^2)\). Thus a positive margin of size
\(c_{\rm p}\nu^2/8\) remains. A contrary sequence of configurations
has a subsequence with bounded ratio or ratio tending to infinity;
both cases exclude its failure. This proves the uniform one-direction
margin on the whole parent. The intersection calculation is the use
of base dimension four.

Where \(\chi=1\) and the period position lies in a core, \(u_0\)
has a fixed positive lower bound. Both directions in \(\mathcal F\)
then have test at least \(c's/\ell\), while the Hessian norm is at most
\((C'+O(C_{\rm p}\nu^2))s/\ell\). Fix sufficiently large
\(C_{\rm p}\) and small \(c_{\rm p}\), then sufficiently small \(\nu\);
this region is prepared at scale \(\ell\).
Retain only whole child cubes contained in it.

To compare mass, tile the affine coordinates by period cells of
side \(\ell\), including a dummy \(z_1\) period. Cells meeting the
inner-parent boundary occupy relative volume \(O(\nu)\).
Discarding child cubes meeting core boundaries loses another
\(O(\nu)\) fraction, because the finite cores have finite boundary
area in a period. On the retained region
\[
 |\nabla v|/s=\sqrt{1+a^2f^2}+O(\nu).
\]
The fixed affine Jacobian cancels against parent volume. Thus
\[
 \frac1{s|P|}\int_{\mathcal R_P}|\nabla v|\,\dd x
 \ge q_{\rm in}C_{\rm core}-C\nu,\qquad
 \frac1{s|P|}\int_P|\nabla h|\,\dd x\le1+C\nu.
\]
Taking both errors below \(\eta_0\) gives a ratio exceeding
\((1+3\eta_0)/(1+\eta_0)>1+\zeta\), with
\(\zeta=\eta_0/(1+\eta_0)>0\).

Finally, \(j\) derivatives of \eqref{eq:modify}, with \(q\) on the
cutoff, cost
\[
 s\ell(\nu S)^{-q}\ell^{-(j-q)}
   =s\ell^{1-j}\nu^q\le s\ell^{1-j}.
\]
The slope and affine frame are frozen constants in this expression.
This proves \eqref{eq:increment}. The order of choices is the
pattern and its small amplitude, the packing and cutoffs,
the preparation constants, and finally \(\nu\).
\end{proof}

\subsection{Iteration with every finite derivative order}

\begin{proof}[Proof of Proposition~\ref{prop:envelope}]
Start from one seed cube of side \(\nu\), prepared at scale \(1\),
compactly inside \(U_0\). Write \(h_0\) for the seed and \(P_0\)
for this cube. Fix \(s_{\min}>0\) with
\( |\nabla_Gh_0|_G\ge s_{\min}\) on \(P_0\), common over
\(G\in\mathcal G\); this follows from the gradient separation on
\(U_0\) and the common metric bounds. Put \(r_{\rm c}=\nu^2\).
At generation \(i\), apply Lemma~\ref{lem:corrugation} on each
selected parent at scale \(S_i=r_{\rm c}^{i-1}\).
The selected children have side \(\nu r_{\rm c}^{i}\).
Supports lie strictly inside parents, so every finite generation
is smooth. The total selected gradient mass \(I_i\) satisfies
\[
 I_i\ge(1+\zeta)I_{i-1},\qquad
 \int_V\sqrt{1+|\nabla h_i|^2}\,\dd x\ge I_i.
\]
The initial \(I_0\) has a common positive lower bound.

Gradients grow by at most \(\beta\) each generation and retain
at least a factor \(\alpha\) on the initial cube. Disjoint supports
within a generation and \eqref{eq:increment} give
\[
 \|D^j(h_i-h_{i-1})\|_\infty
 \le C_j\beta^{i-1}r_{\rm c}^{i(1-j)}.
\]
There is no factor counting cubes. The increments are added, so
their higher derivative constants are not multiplied between
generations. Summing yields one \(C_*\), independent of \(j\), with
\begin{equation}\label{eq:iteration}
 \|h_m\|_{C^j}\le C_j e^{C_*(j+1)m}.
\end{equation}
For example, take \(C_*\) larger than
\(\log\beta+|\log r_{\rm c}|+1\).
At a point last changed at generation \(i\le m\), the retained
Hessian margin is at least
\[
 c\,s_{\min}\alpha^{i-1}/r_{\rm c}^{i-1}
 \ge c\,s_{\min}\alpha^m.
\]
Later modifications vanish there, with all jets vanishing at
support boundaries. Outside the initial cube the seed is unchanged.
Enlarge \(C_*\) to obtain \(\max_eJ_{h_m}(e)\ge c e^{-C_*m}\)
on \(t\ge1/4\).

Choose once \(\theta>0\) with \(C_*\theta\le\delta\), and put
\[
 m=\lfloor\theta\log k\rfloor,\qquad
 \gamma=\theta\log(1+\zeta)>0,\qquad h_k=h_m+c_k.
\]
The derivative and Hessian bounds follow from \eqref{eq:iteration},
and the mass bound follows from
\((1+\zeta)^m\ge(1+\zeta)^{-1}k^\gamma\).
All changes lie in the seed cube above \(t=1/2\).
The constants determining \(C_*,\theta,\gamma\) use the fixed
pattern and common \(C^1\) metric bounds only. A later smooth coordinate
change contributes finite order-dependent chain-rule constants and
lower derivatives of \(h_k\), without changing these exponents.
\end{proof}
